\documentclass[12pt,a4paper]{article}
\usepackage[a4paper,top=15mm,bottom=15mm,left=15mm,right=15mm]{geometry}
\usepackage[T1]{fontenc}
\usepackage{amsmath}
\usepackage{amssymb}

\pagestyle{empty}
\setlength{\parindent}{0pt}
\setlength{\parskip}{6pt}

\begin{document}

\begin{center}
{\Large\bfseries Worked Solutions: An Introduction to Surds}
\end{center}

\section*{Part 1: Square roots on the calculator}

\subsection*{Q1}
The completed table is:

\begin{center}
\renewcommand{\arraystretch}{1.5}
\begin{tabular}{|c|c|c|}
\hline
\textbf{Square root} & \textbf{Calculator answer (4 d.p.)} & \textbf{Is it a surd?}\\
\hline
$\sqrt{16}$ & $4$ & No\\
\hline
$\sqrt{2}$ & $1.4142$ & Yes\\
\hline
$\sqrt{49}$ & $7$ & No\\
\hline
$\sqrt{3}$ & $1.7320$ & Yes\\
\hline
$\sqrt{10}$ & $3.1622$ & Yes\\
\hline
\end{tabular}
\end{center}

A surd is a square root that cannot be written exactly as a whole number or a fraction. So $\sqrt{16}$ and $\sqrt{49}$ are not surds; $\sqrt{2}$, $\sqrt{3}$ and $\sqrt{10}$ are surds.

\section*{Part 2: Surds that are equal}

\subsection*{Q2}
(a) Using the calculator and truncating to 4 decimal places:
\begin{enumerate}
\item $2\sqrt2 = 2.8284$
\item $\sqrt{12} = 3.4641$
\item $\sqrt6 = 2.4494$
\item $\sqrt{18} = 4.2426$
\item $\sqrt{20} = 4.4721$
\item $2\sqrt3 = 3.4641$
\item $\sqrt8 = 2.8284$
\item $3\sqrt5 = 6.7082$
\item $2\sqrt5 = 4.4721$
\item $3\sqrt2 = 4.2426$
\end{enumerate}

(b) The four equal pairs are:
\begin{itemize}
\item $2\sqrt2 = \sqrt8$
\item $\sqrt{12} = 2\sqrt3$
\item $\sqrt{18} = 3\sqrt2$
\item $\sqrt{20} = 2\sqrt5$
\end{itemize}
The unpaired numbers are $\sqrt6$ and $3\sqrt5$.

Explanation: for example, $\sqrt{12} = \sqrt{4\times3} = \sqrt4 \times \sqrt3 = 2\sqrt3$, so $\sqrt{12}$ is equal to $2\sqrt3$.

\subsection*{Q3}
(a)
\begin{align*}
\sqrt{12} &= \sqrt{4\times3} \\
&= \sqrt4 \times \sqrt3 \\
&= 2\sqrt3.
\end{align*}

(b)
\begin{align*}
\sqrt{50} &= \sqrt{25\times2} \\
&= \sqrt{25} \times \sqrt2 \\
&= 5\sqrt2.
\end{align*}

(c) $\sqrt{24} = \sqrt{4\times6} = 2\sqrt6$.

(d) $\sqrt{28} = \sqrt{4\times7} = 2\sqrt7$.

(e) $\sqrt{45} = \sqrt{9\times5} = 3\sqrt5$.

(f) $\sqrt{63} = \sqrt{9\times7} = 3\sqrt7$.

(g) $\sqrt{75} = \sqrt{25\times3} = 5\sqrt3$.

(h) $\sqrt{98} = \sqrt{49\times2} = 7\sqrt2$.

\section*{Part 3: Adding and subtracting surds}

\subsection*{Q4}
(a) $2x+4x = (2+4)x = 6x$.

(b) $5y-2y = (5-2)y = 3y$.

(c) $a+3a = 1a+3a = 4a$.

(d) $6m-m = 6m-1m = 5m$.

(e) $4x+2y+x = (4+1)x+2y = 5x+2y$.

(f) $x+y$ cannot be simplified because $x$ and $y$ are not like terms.

\subsection*{Q5}
(a) $3\sqrt5+2\sqrt5 = (3+2)\sqrt5 = 5\sqrt5$.

(b) $9\sqrt7-4\sqrt7 = (9-4)\sqrt7 = 5\sqrt7$.

(c) $4\sqrt2+\sqrt2 = 4\sqrt2+1\sqrt2 = 5\sqrt2$.

(d) $3\sqrt6+5\sqrt6 = 8\sqrt6$.

(e) $8\sqrt{11}-3\sqrt{11} = 5\sqrt{11}$.

(f) $6\sqrt5-\sqrt5 = 6\sqrt5-1\sqrt5 = 5\sqrt5$.

(g) $2\sqrt3+3\sqrt3+\sqrt3 = (2+3+1)\sqrt3 = 6\sqrt3$.

(h) $5\sqrt2+3\sqrt7+2\sqrt2 = (5+2)\sqrt2+3\sqrt7 = 7\sqrt2+3\sqrt7$.

(i) $\sqrt2+\sqrt3$ cannot be simplified because $\sqrt2$ and $\sqrt3$ are not like surds. It is not $\sqrt5$.

(j) $4\sqrt6-4\sqrt6 = 0$.

\subsection*{Q6}
(a) The rectangle has two sides of length $3\sqrt2$ and two sides of length $\sqrt2$. So
\[
\text{Perimeter} = 3\sqrt2+\sqrt2+3\sqrt2+\sqrt2 = (3+1+3+1)\sqrt2 = 8\sqrt2 \text{ cm}.
\]

(b) The triangle has sides $2\sqrt5$, $3\sqrt5$ and $4\sqrt5$. So
\[
\text{Perimeter} = 2\sqrt5+3\sqrt5+4\sqrt5 = (2+3+4)\sqrt5 = 9\sqrt5 \text{ cm}.
\]

\subsection*{Q7}
(a) $\sqrt8+\sqrt2 = \sqrt{4\times2}+\sqrt2 = 2\sqrt2+\sqrt2 = 3\sqrt2$.

(b) $\sqrt{12}+\sqrt3 = \sqrt{4\times3}+\sqrt3 = 2\sqrt3+\sqrt3 = 3\sqrt3$.

(c) $\sqrt{18}-\sqrt2 = \sqrt{9\times2}-\sqrt2 = 3\sqrt2-\sqrt2 = 2\sqrt2$.

(d) $\sqrt{45}+\sqrt{20} = \sqrt{9\times5}+\sqrt{4\times5} = 3\sqrt5+2\sqrt5 = 5\sqrt5$.

\section*{Part 4: Multiplying and dividing surds}

\subsection*{Q8}
(a) $\sqrt2\times\sqrt7 = \sqrt{2\times7} = \sqrt{14}$.

(b) $3\sqrt2\times5\sqrt7 = (3\times5)\sqrt{2\times7} = 15\sqrt{14}$.

(c) $\sqrt3\times\sqrt7 = \sqrt{21}$.

(d) $\sqrt6\times\sqrt6 = 6$.

(e) $\sqrt2\times\sqrt{18} = \sqrt{36} = 6$.

(f) $\sqrt5\times\sqrt{20} = \sqrt{100} = 10$.

(g) $2\sqrt5\times3\sqrt2 = (2\times3)\sqrt{5\times2} = 6\sqrt{10}$.

(h) $4\sqrt3\times2\sqrt7 = (4\times2)\sqrt{3\times7} = 8\sqrt{21}$.

(i) $5\sqrt2\times\sqrt3 = 5\sqrt{6}$.

(j) $3\sqrt5\times3\sqrt5 = (3\times3)(\sqrt5\times\sqrt5) = 9\times5 = 45$.

\subsection*{Q9}
(a) $\sqrt{21}\div\sqrt7 = \sqrt{21\div7} = \sqrt3$.

(b) $10\sqrt{15}\div5\sqrt5 = (10\div5)\sqrt{15\div5} = 2\sqrt3$.

(c) $\sqrt{30}\div\sqrt6 = \sqrt{30\div6} = \sqrt5$.

(d) $\dfrac{\sqrt{22}}{\sqrt2} = \sqrt{\dfrac{22}{2}} = \sqrt{11}$.

(e) $\dfrac{\sqrt{12}}{\sqrt3} = \sqrt{\dfrac{12}{3}} = \sqrt4 = 2$.

(f) $\sqrt{63}\div\sqrt7 = \sqrt{63\div7} = \sqrt9 = 3$.

(g) $8\sqrt6\div4\sqrt2 = (8\div4)\sqrt{6\div2} = 2\sqrt3$.

(h) $12\sqrt{10}\div3\sqrt2 = (12\div3)\sqrt{10\div2} = 4\sqrt5$.

\subsection*{Q10}
(a) Area $= \sqrt2\times\sqrt8 = \sqrt{2\times8} = \sqrt{16} = 4$ cm$^2$.

(b) Area $= 5\sqrt2\times2\sqrt3 = (5\times2)\sqrt{2\times3} = 10\sqrt6$ cm$^2$.

(c) Length $= \dfrac{\text{Area}}{\text{width}} = \dfrac{\sqrt{35}}{\sqrt5} = \sqrt{\dfrac{35}{5}} = \sqrt7$ cm.

\section*{Part 5: Many ways to make a surd}

\subsection*{Q11}
For each part, one valid completion is given.

(a) Middle $5\sqrt2$:
\begin{itemize}
\item Addition: $\sqrt2+4\sqrt2 = 5\sqrt2$.
\item Subtraction: $8\sqrt2-3\sqrt2 = 5\sqrt2$.
\item Multiplication: $\sqrt{25}\times\sqrt2 = 5\sqrt2$.
\item Division: $10\sqrt6\div2\sqrt3 = 5\sqrt2$.
\item Single square root: $\sqrt{50} = 5\sqrt2$.
\end{itemize}

(b) Middle $3\sqrt5$:
\begin{itemize}
\item Addition: $\sqrt5+2\sqrt5 = 3\sqrt5$ (given).
\item Subtraction: $4\sqrt5-\sqrt5 = 3\sqrt5$.
\item Multiplication: $\sqrt9\times\sqrt5 = 3\sqrt5$.
\item Division: $6\sqrt{10}\div2\sqrt2 = 3\sqrt5$.
\item Single square root: $\sqrt{45} = 3\sqrt5$ (given).
\end{itemize}

(c) Middle $6\sqrt2$:
\begin{itemize}
\item Addition: $2\sqrt2+4\sqrt2 = 6\sqrt2$.
\item Subtraction: $7\sqrt2-\sqrt2 = 6\sqrt2$.
\item Multiplication: $\sqrt{36}\times\sqrt2 = 6\sqrt2$.
\item Division: $12\sqrt6\div2\sqrt3 = 6\sqrt2$.
\item Single square root: $\sqrt{72} = 6\sqrt2$.
\end{itemize}

(d) Middle $2\sqrt7$:
\begin{itemize}
\item Addition: $\sqrt7+\sqrt7 = 2\sqrt7$.
\item Subtraction: $5\sqrt7-3\sqrt7 = 2\sqrt7$.
\item Multiplication: $\sqrt4\times\sqrt7 = 2\sqrt7$.
\item Division: $6\sqrt{14}\div3\sqrt2 = 2\sqrt7$.
\item Single square root: $\sqrt{28} = 2\sqrt7$.
\end{itemize}

(e) One possible choice: middle $3\sqrt2$.
\begin{itemize}
\item Addition: $\sqrt2+2\sqrt2 = 3\sqrt2$.
\item Subtraction: $5\sqrt2-2\sqrt2 = 3\sqrt2$.
\item Multiplication: $\sqrt9\times\sqrt2 = 3\sqrt2$.
\item Division: $9\sqrt6\div3\sqrt3 = 3\sqrt2$.
\item Single square root: $\sqrt{18} = 3\sqrt2$.
\end{itemize}

In every case, check the box by simplifying the calculation to see that it gives the middle surd.

\end{document}